\section{Fair sector search and exact Euler-LP crushing}
\label{sec:euler-lp-crushing}

The preceding complete crushing backend can spend its entire finite allowance
in one sector before considering a later easy positive surface.  It also
enumerates Q corners even when a short linear feasibility certificate could
settle canonical positivity.  This section develops two improvements: a
deterministic restart schedule and an exact Euler-LP route with independent
Farkas replay.  These are changes to discovery and its evidence, not new
topological premises.  The diagram-bound crushing theorem and strict
tetrahedron descent remain those of Section~\ref{sec:normal-crushing-search}.

The LP implementation selectively adapts the solver and Euler certificate
ideas preserved in report~67, discussed in
Section~\ref{sec:completion-theorems}.  Its archived files remain unchanged.
The maintained version adds current source checks, cancellation handling,
binary certificate transport, support extraction and diagram-level use.
Neither the complete sector family nor the exact simplex solver currently
has a general quasi-polynomial implementation bound.

\subsection{Canonical Euler characteristic as finitely many linear objectives}

Fix a compatible allowed quadrilateral set of size $k\leq t$.  Source
matching equations give a nonnegative Q cone
\[
 Q=\{q\in\mathbb R^k:Cq=0,\ q\geq0\}.
\]
After contracting zero-labelled triangle equations, let $G_v$ be the
remaining corner classes at vertex $v$, with linear triangle potentials
$\ell_{vc}$.  Constant groups are absorbed into the linear part.  For the
nonconstant groups, canonical Euler characteristic is
\[
 F(q)=L(q)-\sum_v w_v\min_{c\in G_v}\ell_{vc}(q),
 \qquad w_v=\chi(\operatorname{link}v)\in\{1,2\}.
\]
The source validator supplies the finite compact manifold premise behind
these positive weights.  Choosing one corner class independently at every
active vertex gives the identity
\[
 F(q)=\max_{\alpha\in\prod_vG_v} c_\alpha^{\mathsf T}q,
 \qquad c_\alpha=L-\sum_v w_v\ell_{v,\alpha_v}.
\]
Consequently $F>0$ somewhere on $Q$ if and only if one homogeneous rational
LP is feasible:
\[
 Cq=0,\qquad q\geq0,\qquad c_\alpha^{\mathsf T}q\geq1.
\]
No inequalities asserting that the chosen anchors are actual minima are
needed.  The implication $c_\alpha^{\mathsf T}q>0\Rightarrow F(q)>0$
holds for every anchor tuple; choosing actual minima proves the converse.

The producer solves the normalized problem $Cq=0$, $q\geq0$,
$\sum q_j\leq1$, stopping at the first positive basic point.  A feasible
slack basis for $Cq\leq0$, $-Cq\leq0$ and $\sum q_j\leq1$ avoids a separate
phase-I solve.  Exact Bland pivots terminate, but their number has no
polynomial bound.  All rational numerators and denominators are carried
exactly; floating positivity tolerances are not used.

\subsection{Independent negative certificates and an envelope shortcut}

For a nonpositive anchor objective the solver returns a rational multiplier
$y_\alpha$ satisfying
\[
 C^{\mathsf T}y_\alpha\geq c_\alpha
\]
componentwise.  If $q\geq0$ and $Cq=0$, this inequality implies
$c_\alpha^{\mathsf T}q\leq y_\alpha^{\mathsf T}Cq=0$.
Checking all anchor tuples therefore excludes positive canonical Euler
throughout the supplied sector.  Even with no active groups the Cartesian
product contains its single empty tuple; an empty list of anchor proofs
does not discharge that obligation.

The checker reconstructs triangle components, potentials, cycle equations
and Euler coefficients from the original source matching rows.  It imports
no LP solver, sector constructor or ray enumerator.  It requires every
anchor in canonical order with the exact rational inequality, or one
componentwise envelope certificate.  A valid packet proves precisely the
source and allowed sector named in it.

The envelope objective
\[
 e_j=L_j-\sum_v w_v\min_{c\in G_v}\ell_{vc,j}
\]
dominates every $c_\alpha$ componentwise.  Thus one multiplier with
$C^{\mathsf T}y\geq e$ excludes the whole sector.  A positive envelope
point by itself is inconclusive: for example, on $q=(s,s)$ the two
objectives $(1,-2)$ and $(-2,1)$ both have value $-s$, while their envelope
$(1,1)$ has value $2s$.

The maintained producer nevertheless evaluates the actual $F(q)$ at a
positive envelope point.  If this value is positive, it constructs the
canonical lift and checks its source-derived Euler characteristic before
returning a witness.  Otherwise it continues the complete anchor search.
This additional check can avoid anchor enumeration without mistaking the
linear relaxation's objective for the surface's Euler value.

Source-row forcing can delete Q coordinates before the LP.  The new
reduction wrapper retains the original sector, the full matching implication
proof, and the Euler packet on exactly the retained types.  Its independent
checker proves every forced zero from original matching equations before
replaying the smaller sector.  A proof about an arbitrary smaller sector
cannot stand in for the original negative query.  Nested reduction wrappers,
altered retained masks and mismatched source digests are rejected.

\subsection{Turning any canonical positive witness into a crushable ray}

A positive Euler certificate intentionally proves a canonical vector, not
connectedness or extremality.  This permits future exact feasibility
methods to return a nonbasic point.  The maintained extractor supplies the
remaining ray contract without expanding normal discs.

Let $x$ be the independently validated positive canonical standard vector,
$S=\operatorname{supp}x$, and $M$ the original matching matrix.  Prepare the
kernel only of $M_{[:,S]}$, retain the exact original Euler coefficients on
$S$, and apply the support descent of the preceding reusable-kernel section.
After primitive rescaling, zero-expand the resulting ray into the original
$7t$ coordinates.

\begin{proposition}[Positive support extraction for crushing]
This procedure returns a primitive admissible positive full matching-cone
extreme ray with a quadrilateral and support contained in $S$.
\end{proposition}
\begin{proof}
The nonnegative cone restricted to $S$ is a coordinate face of the full
matching cone, and its extreme rays are extreme in that full cone.  The
descent preserves every initial zero and never decreases the normalized
positive Euler objective.  The retained support contains no incompatible
quad pair.  At every vertex at least one triangle coordinate is initially
zero, because $x$ is canonical; these zeros persist.  A nonzero triangle-only
matching vector would be a nonnegative sum of vertex links and would violate
one of those zeros.  The output therefore has a quadrilateral.  Primitive
integral scaling and full-source extremality give connectedness by the
lemma in Section~\ref{sec:normal-crushing-search}.  The crushing verifier
independently rechecks the complete original source and ray contract.
\end{proof}

Preparing the positive support instead of all $7t$ columns is especially
useful for these one-shot queries.  The extractor still retains every
original matching equation; it does not impose an unrelated linear slice
that could make a disconnected surface appear extreme.

\subsection{Growing prefixes, fair restarts and shared allowances}

Round $r$ visits the first $\min(16\cdot2^r,3^t)$ assignments in the selected
deterministic order, allowing $q_0 2^r$ candidates for each unfinished
sector, with $q_0=64$ by default.  Completed negative sectors are skipped.
An unfinished corner job retains its prepared kernel and matching-reduction
proof, but restarts its deterministic enumerator; a partial transcript is
discarded.  The growing prefix gives later sectors an early opportunity,
while growing the quota revisits expensive early sectors before traversing
the entire exponential family.  The sequential policy remains available
as \texttt{sector_quantum=None}.

For a deterministic corner job that completes or finds its first positive
after $W>0$ candidates, all aborted attempt quotas are less than $W$.
They form a geometric sequence whose sum is less than $2W$.  Adding the
final $W$ candidates bounds that job's repeated candidate work by $3W$.
Zero-candidate jobs complete on their first visit.  Source preparation is
cached for unfinished jobs; enumerator initialization, prefix traversal
and certificate processing remain additional work.  Cached kernels, the
explored assignment prefix and negative certificates can still occupy
exponential memory.  Fairness supplies no small global sector count.

The Euler-LP option uses the same growing schedule with separate per-attempt
pivot and anchor quotas.  It rebuilds its LP state on retry and charges every
repeated pivot and anchor.  Source matching preparation remains cached in
the enclosing job.  Since a quota can stop either resource, the clean
geometric bound concerns the sum of pivot and anchor counts, rather than
a separate constant-factor guarantee for each count.  All unbounded jobs
are finite: anchor products are finite and exact Bland pivoting terminates.
Both complete modes therefore retain the strict-descent completeness proof.

The public crushing API keeps corner search as its default and exposes
\texttt{sector_method='lp'} with \texttt{lp_strategy='envelope'} or
\texttt{'anchors'}.  Global sector starts and crushing steps are shared
across successors.  Corner bases, LP pivots and LP anchors have explicit,
separate allowances; only the selected method's work counts are consumed.
Reaching an applicable allowance returns an authenticated partial chain,
not a negative sector proof.  Caller exceptions remain distinct from those
local limits and propagate through nested replay unchanged.

For a final knottedness packet the global verifier still requires exactly
all $3^t$ maximal assignments in canonical order.  Each can carry a complete
Q-corner exhaustion, a native Euler-LP proof, or a source-proven matching
reduction followed by an Euler-LP proof.  The checker independently verifies
each packet's full original assignment.  Neither the traversal order nor
mixed proof methods permit a missing sector or a partial negative claim.

\subsection{What a fixed-sector parameter bound would prove}

Let $g$ be the number of nonconstant corner groups and $m_v=|G_v|$.  Their
total number of classes is at most $4t$.  A sharper support bound follows
from the triangle incidence graph: allowing $k$ quad coordinates removes
at most $4k$ zero-labelled triangle adjacency edges.  Since the original
corner graph at each vertex is connected, contraction leaves
$\sum_v(m_v-1)\leq4k$.  Hence, when $g>0$, the number of anchor objectives
satisfies
\[
 A=\prod_v m_v\leq
 \min\{(4t/g)^g,\ (1+4k/g)^g\}.
\]
For $g=0$ the count is one.  Identifying potentials equal on $\ker C$ can
only reduce this bound, but the current certificate format requires its
explicit source-derived corner groups and does not silently omit anchors.

The producer retains at most 256 cached duals, always including the zero
dual, so coordinatewise dominance tests add only polynomial work per
anchor.  An unbounded all-previous-duals scan would instead permit quadratic
work in $A$; no such scan is needed for correctness.

With a polynomial exact LP policy this would give $A\operatorname{poly}(t)$
fixed-sector discovery and replay, followed by polynomial support extraction
and crushing.  One active vertex needs at most $\min(4t,4k+1)$ anchor LPs,
plus at most one preliminary envelope LP;
$g=O(\log t)$ gives a quasi-polynomial fixed-sector bound.  These are
conditional time bounds, not bounds for the maintained Bland solver.
Moreover, the full $3^t$ compatible sector family remains even with a
polynomial solver.  A uniform small sector/backdoor family or a suitable
source decomposition on every crushing successor is still required for
the requested general quasi-polynomial recognizer.
